\chapter{Podsumowanie}
\label{cha:podsumowanie}

% \section{Realizacja celu pracy}
% \label{sec:realizacja_celu}

Celem pracy było opracowanie systemu wyznaczającego ruch robota na podstawie danych z~dwóch kamer zdarzeniowych, przeznaczonego do przetwarzania zbiorów EvSLAM i~M3ED. Cel ten został zrealizowany. Powstał system zawierający moduły odpowiedzialne za synchronizację i~przygotowanie zdarzeń, śledzenie punktów, geometrię stereo oraz estymację trajektorii lewej kamery. Dla EvSLAM wyznaczana jest również wymagana przez konkurs prędkość liniowa.

Odometria wizyjna została rozszerzona o~rzadką mapę, lokalizację względem jej punktów, relokalizację, domykanie pętli i~optymalizację grafu pozy. System przetworzył wszystkie oceniane sekwencje EvSLAM oraz wybrane pełne przebiegi M3ED. Działa on na gotowych, zebranych wcześniej sekwencjach i~obecnie nie jest przystosowany do pracy w~czasie rzeczywistym na robocie mobilnym. Wykorzystuje IMU, lecz używa wyłącznie żyroskopu, dlatego nie stanowi pełnego estymatora wizyjno-inercyjnego. Prędkość oblicza z~gotowej trajektorii położenia, co jest prostym rozwiązaniem stosowanym na potrzeby uzyskania wyników wymaganych przez EvSLAM.

\section{Kierunki dalszych prac}
\label{sec:dalsze_prace}

Najważniejsze kierunki rozwoju, w~kolejności wynikającej z~wykazanych ograniczeń, obejmują:

\begin{enumerate}
    \item utworzenie pełnego estymatora wizyjno-inercyjnego, który wspólnie wyznaczałby pozę, prędkość oraz błędy systematyczne IMU z~wykorzystaniem preintegracji pomiarów~\cite{forster2015preintegration} -- mogłoby to ograniczyć dryf rotacji, będący jednym z~najważniejszych ograniczeń widocznych w~wynikach dla sekwencji \texttt{falcon\_outdoor\_day\_fast\_flight\_2};
    \item uwzględnienie niepewności estymacji głębi podczas wyboru punktów do estymacji pozy;
    \item wprowadzenie adaptacyjnego doboru czasu akumulacji zależnego od dynamiki sceny, który zwiększyłby odporność systemu na zmiany prędkości ruchu robota i~pozwoliłby skutecznie estymować pozę z~większej liczby klatek zdarzeniowych;
    \item dodanie lokalnej optymalizacji wiązki (ang. \textit{bundle adjustment}), stosowanej w~bardziej rozbudowanych systemach SLAM, między innymi w~ORB-SLAM2~\cite{murartal2017orbslam2}, oraz zastosowanie bardziej wydajnego optymalizatora grafowego dla większych map~\cite{grisetti2010graph};
    \item ocena złożoności obliczeniowej i~przyspieszenie najbardziej kosztownych etapów, a~następnie sprawdzenie systemu na bieżącym strumieniu danych.
\end{enumerate}
